\documentclass[conference]{IEEEtran}

\usepackage{cite}
\usepackage{amsmath,amssymb,amsfonts}
\usepackage{graphicx}
\usepackage{booktabs}
\usepackage{array}
\usepackage{url}
\usepackage{xcolor}

\title{Notes on Linear Regression}

\author{
\IEEEauthorblockN{Luis Mendez}
}

\begin{document}
\maketitle

\textbf{Question:} What is linear regression doing, and how does a
neural network learn it?

\bigskip

\textbf{Answer:}

\bigskip

\textbf{1. The Model}

Linear regression fits a straight line through data: given input $x$, it
predicts

\[
  \hat{y} = a x + b,
\]

where $a$ is the \textbf{slope} and $b$ is the \textbf{intercept}. With
$D$ input features this generalizes to $\hat{y} = \mathbf{w}^\top
\mathbf{x} + b$ --- exactly a single neuron with \emph{no} activation
function (or, equivalently, a linear $\phi$). Framing it this way, a
dense layer of one unit with linear activation \emph{is} linear
regression.

\bigskip

\textbf{2. Loss: Mean Squared Error}

Training picks $a, b$ to minimize the squared distance between
predictions and targets, averaged over $N$ examples:

\[
  C(a,b) = \frac{1}{N}\sum_{i=1}^{N} \left(y_i - (a x_i + b)\right)^2.
\]

Squaring penalizes large errors more than small ones and keeps $C$
smooth and differentiable everywhere, which is what makes both routes
below possible.

\bigskip

\textbf{3. Two Ways to Solve It}

\begin{itemize}
    \item \textbf{Closed form (OLS):} because $C$ is quadratic and
          convex in $a,b$, setting $\partial C/\partial a =
          \partial C/\partial b = 0$ and solving directly gives

\[
  a = \frac{\displaystyle\sum x_i y_i - \bar{y}\sum x_i}
           {\displaystyle\sum x_i^2 - \bar{x}\sum x_i},
  \qquad
  b = \bar{y} - a\bar{x}.
\]

          No iteration needed --- one pass over the data.
    \item \textbf{Gradient descent:} treat $a,b$ as the weight and bias
          of a one-neuron network and train it the same way as any
          other neuron (see \textit{Notes on Neurons}): repeatedly
          nudge $a,b$ opposite $\nabla C$ until the loss stops
          improving. Slower and approximate, but the only option once
          the model has more layers or nonlinearity, where no closed
          form exists.
\end{itemize}

Both routes should agree, up to numerical precision, on a purely linear
problem --- a good sanity check after training.

\bigskip

\textbf{4. Linearizing an Exponential Trend}

Not every relationship is linear on its own scale, but some become
linear after a transform. Moore's Law states that transistor count
roughly doubles every two years, i.e. $y \approx y_0 \cdot 2^{x/2}$, an
exponential in the raw data. Taking the log of both sides,

\[
  \log y = \log y_0 + \frac{\log 2}{2} x,
\]

which is linear in $x$. So $\log y$ vs.\ $x$ can be fit with ordinary
linear regression, and the fitted slope $a$ converts back into a
doubling time via $a = \tfrac{\log 2}{2\,\text{years}} \Rightarrow
\text{doubling time} = \log 2 / a$. Centering $x$ (subtracting its mean)
before fitting does not change $a$; it only keeps $b$ interpretable and
improves numerical conditioning for gradient descent.

\bigskip

\textbf{5. Keras Example}

\begin{verbatim}
model = tf.keras.Sequential([
    tf.keras.layers.Input(shape=(1,)),
    tf.keras.layers.Dense(1),
])

model.compile(
    optimizer=tf.keras.optimizers.SGD(0.001, 0.9),
    loss='mse',
)

r = model.fit(X, Y, epochs=100)
\end{verbatim}

A single \texttt{Dense(1)} layer with the default (linear) activation is
the whole model: one weight, one bias. The learned weight and bias
recovered from \texttt{model.layers[0].get\_weights()} should match the
closed-form $a,b$ computed directly from the data.

\bigskip

\textbf{6. Saving and Reloading}

\begin{verbatim}
model.save('linearregression.h5')

loaded = tf.keras.models.load_model('linearregression.h5')
np.allclose(model.predict(X), loaded.predict(X))
\end{verbatim}

Saving stores the architecture, learned weights, and optimizer state;
reloading should reproduce identical predictions, which is a useful
check that nothing was lost in the round trip.

\bigskip

\textbf{7. Putting It Together}

\begin{itemize}
    \item Linear regression = one neuron, linear activation, MSE loss.
    \item With one input feature it can be solved exactly (OLS) or
          learned iteratively (gradient descent); both should agree.
    \item A log transform turns an exponential trend (e.g.\ Moore's Law)
          into a linear one, so the same tools apply.
    \item The same neuron-and-loss machinery scales up to deep networks
          once activations become nonlinear and the closed form
          disappears.
\end{itemize}

\end{document}
